\documentclass[a4paper,11pt]{article}
\usepackage[utf8]{inputenc}
\usepackage[T1]{fontenc}
\usepackage[spanish]{babel}
\usepackage[margin=2.5cm]{geometry}
\usepackage{enumitem}
\usepackage{xcolor}
\usepackage{titlesec}
\usepackage{parskip}

% Colores
\definecolor{unidadcolor}{RGB}{30, 60, 120}
\definecolor{practicacolor}{RGB}{0, 100, 0}

% Formato de títulos
\titleformat{\section}
  {\Large\bfseries\color{unidadcolor}}
  {\thesection.}{0.5em}{}
\titleformat{\subsection}
  {\large\bfseries}
  {\thesubsection}{0.5em}{}

\begin{document}

\begin{center}
    {\LARGE \textbf{Métodos de Análisis de Datos 1}}\\[0.5em]
    {\large Programa 2026 --- Primer Cuatrimestre}\\[0.3em]
    {\normalsize Universidad Nacional del Sur}\\[0.3em]
    {\normalsize Departamento de Matemática}\\[0.3em]
    {\normalsize Profesor: José Bavio}
\end{center}

\vspace{1em}

\noindent\textbf{Carga horaria:} 120 horas (30 clases de 2 horas, lunes y miércoles)\\
\textbf{Modalidad:} Presencial\\
\textbf{Herramientas:} Python, R, VS Code, Jupyter Notebooks

\vspace{1em}

\section*{Objetivos}

Al finalizar el curso, los estudiantes serán capaces de:
\begin{itemize}[leftmargin=*]
    \item Aplicar el ciclo completo de análisis de datos: obtención, limpieza, exploración, modelado y comunicación
    \item Utilizar Python y R para tareas de análisis de datos
    \item Construir y evaluar modelos estadísticos y de machine learning
    \item Comunicar hallazgos mediante visualizaciones y reportes reproducibles
    \item Trabajar con datos reales de diversas fuentes
\end{itemize}

\vspace{1em}
\setcounter{section}{-1}

\section{Unidad 0: Introducción al Análisis de Datos}

\subsection{Contenidos teóricos}
\begin{itemize}[leftmargin=*]
    \item El proceso de análisis de datos: ciclo obtención → limpieza → exploración → modelado → predicción → comunicación
    \item Tipos de análisis: descriptivo, exploratorio, inferencial, predictivo, causal
    \item Roles en ciencia de datos: analista, científico de datos, ingeniero de datos
    \item Configuración del entorno de trabajo: instalación y configuración de VS Code, Python (Anaconda), R
    \item Introducción a Jupyter notebooks y R Markdown
    \item Sintaxis básica de Python y R: variables, tipos de datos, estructuras de control
    \item Buenas prácticas: código limpio, comentarios, control de versiones con Git
    \item Reproducibilidad en análisis de datos
    \item Caso de estudio completo: de los datos crudos al reporte final
\end{itemize}

\subsection{Práctica}
\begin{itemize}[leftmargin=*, labelindent=0pt, label={\textcolor{practicacolor}{$\blacktriangleright$}}]
    \item \textbf{Práctica 0.1:} Instalación y configuración del entorno (Python, R, VS Code, extensiones)
    \item \textbf{Práctica 0.2:} Primeros notebooks en Python y R: operaciones básicas, importar librerías
    \item \textbf{Práctica 0.3:} Análisis exploratorio de un dataset sencillo (iris, mtcars) siguiendo el ciclo completo
\end{itemize}

\section{Unidad 1: Obtención y Preparación de Datos}

\subsection{Contenidos teóricos}
\begin{itemize}[leftmargin=*]
    \item Fuentes de datos: archivos planos (CSV, TSV, Excel), bases de datos (SQL), APIs REST, web scraping
    \item Carga de datos en Python: pandas (read\_csv, read\_excel, read\_sql)
    \item Carga de datos en R: readr, readxl, DBI
    \item Exploración inicial: dimensiones, tipos de datos, primeras filas, resumen estadístico
    \item Calidad de datos: completitud, consistencia, precisión, actualidad
    \item Limpieza de datos:
    \begin{itemize}
        \item Detección y tratamiento de valores faltantes (NA, NaN, NULL)
        \item Estrategias: eliminación, imputación (media, mediana, moda, KNN, regresión)
        \item Duplicados: detección y eliminación
        \item Inconsistencias: valores fuera de rango, formatos incorrectos
    \end{itemize}
    \item Transformación de datos:
    \begin{itemize}
        \item Conversión de tipos de datos (casting)
        \item Creación de variables derivadas (feature engineering básico)
        \item Codificación de variables categóricas: one-hot encoding, label encoding
        \item Normalización y estandarización
        \item Discretización de variables continuas (binning)
    \end{itemize}
    \item Integración de datos: joins, merges, concatenación
    \item Reshape: pivot, melt, stack, unstack
    \item Tidy data: principios de Hadley Wickham
\end{itemize}

\subsection{Práctica}
\begin{itemize}[leftmargin=*, labelindent=0pt, label={\textcolor{practicacolor}{$\blacktriangleright$}}]
    \item \textbf{Práctica 1.1:} Carga de datos desde CSV y Excel en Python y R
    \item \textbf{Práctica 1.2:} Conexión a base de datos SQL y extracción de datos
    \item \textbf{Práctica 1.3:} Web scraping básico con BeautifulSoup/rvest
    \item \textbf{Práctica 1.4:} Detección y tratamiento de valores faltantes (múltiples estrategias)
    \item \textbf{Práctica 1.5:} Limpieza de un dataset "sucio": duplicados, inconsistencias, formatos
    \item \textbf{Práctica 1.6:} Transformaciones: creación de variables, codificación, normalización
    \item \textbf{Práctica 1.7:} Integración de múltiples fuentes de datos (joins y merges)
\end{itemize}

\section{Unidad 2: Exploración y Visualización de Datos}

\subsection{Contenidos teóricos}
\begin{itemize}[leftmargin=*]
    \item Estadística descriptiva univariada:
    \begin{itemize}
        \item Medidas de tendencia central: media, mediana, moda
        \item Medidas de dispersión: varianza, desviación estándar, rango intercuartílico
        \item Medidas de forma: asimetría (skewness), curtosis
        \item Cuantiles y percentiles
    \end{itemize}
    \item Estadística descriptiva bivariada:
    \begin{itemize}
        \item Tablas de contingencia
        \item Covarianza y correlación (Pearson, Spearman, Kendall)
        \item Interpretación de correlaciones
    \end{itemize}
    \item Análisis de distribuciones:
    \begin{itemize}
        \item Distribuciones empíricas vs teóricas
        \item Test de normalidad (Shapiro-Wilk, Kolmogorov-Smirnov)
        \item Q-Q plots
    \end{itemize}
    \item Detección de outliers:
    \begin{itemize}
        \item Métodos univariados: IQR, Z-score
        \item Métodos multivariados: Mahalanobis distance, isolation forest
    \end{itemize}
    \item Visualizaciones en Python:
    \begin{itemize}
        \item matplotlib: anatomía de un gráfico, subplots
        \item seaborn: gráficos estadísticos de alto nivel
        \item plotly: gráficos interactivos
        \item Tipos: histogramas, boxplots, scatter plots, line plots, heatmaps, pair plots
    \end{itemize}
    \item Visualizaciones en R:
    \begin{itemize}
        \item ggplot2: gramática de gráficos
        \item Capas, aesthetic mappings, geometrías, facets
        \item Personalización: themes, scales, labels
    \end{itemize}
    \item Principios de visualización efectiva:
    \begin{itemize}
        \item Elegir el gráfico apropiado según el tipo de datos y mensaje
        \item Uso efectivo del color, tamaño, forma
        \item Evitar distorsiones y gráficos engañosos
        \item Narrativa visual: contar una historia con los datos
    \end{itemize}
    \item Exploración de datos multivariados: PCA (introducción conceptual)
\end{itemize}

\subsection{Práctica}
\begin{itemize}[leftmargin=*, labelindent=0pt, label={\textcolor{practicacolor}{$\blacktriangleright$}}]
    \item \textbf{Práctica 2.1:} Estadística descriptiva univariada en Python y R
    \item \textbf{Práctica 2.2:} Análisis bivariado: correlaciones y tablas de contingencia
    \item \textbf{Práctica 2.3:} Visualizaciones básicas con matplotlib y ggplot2
    \item \textbf{Práctica 2.4:} Gráficos estadísticos con seaborn
    \item \textbf{Práctica 2.5:} Análisis de distribuciones y tests de normalidad
    \item \textbf{Práctica 2.6:} Detección de outliers (múltiples métodos)
    \item \textbf{Práctica 2.7:} Visualizaciones interactivas con plotly
    \item \textbf{Práctica 2.8:} Exploración multivariada: pair plots, heatmaps de correlación
    \item \textbf{Práctica 2.9:} Proyecto: análisis exploratorio completo de un dataset real
\end{itemize}

\section{Unidad 3: Modelado Estadístico}

\subsection{Contenidos teóricos}
\begin{itemize}[leftmargin=*]
    \item Correlación vs causalidad:
    \begin{itemize}
        \item Variables confundidoras
        \item Experimentos vs estudios observacionales
        \item Introducción a inferencia causal (conceptual)
    \end{itemize}
    \item Regresión lineal simple:
    \begin{itemize}
        \item Formulación del modelo: $y = \beta_0 + \beta_1 x + \epsilon$
        \item Método de mínimos cuadrados ordinarios (OLS)
        \item Interpretación de coeficientes
        \item Supuestos del modelo lineal: linealidad, normalidad, homocedasticidad, independencia
    \end{itemize}
    \item Regresión lineal múltiple:
    \begin{itemize}
        \item Extensión a múltiples predictores
        \item Interpretación de coeficientes parciales
        \item Multicolinealidad: detección (VIF) y tratamiento
        \item Selección de variables: forward, backward, stepwise
    \end{itemize}
    \item Evaluación de modelos de regresión:
    \begin{itemize}
        \item Coeficiente de determinación: $R^2$ y $R^2$ ajustado
        \item Error cuadrático medio: MSE, RMSE
        \item Error absoluto medio: MAE
        \item Análisis de residuos: gráficos diagnósticos
        \item Tests de hipótesis sobre coeficientes
        \item Intervalos de confianza y predicción
    \end{itemize}
    \item Regularización:
    \begin{itemize}
        \item Motivación: overfitting vs underfitting, bias-variance tradeoff
        \item Ridge regression (L2): penalización de la magnitud de coeficientes
        \item LASSO (L1): penalización y selección de variables
        \item Elastic Net: combinación de L1 y L2
        \item Selección de hiperparámetros: cross-validation
    \end{itemize}
    \item Regresión logística:
    \begin{itemize}
        \item Clasificación binaria
        \item Función logística (sigmoide)
        \item Interpretación de odds ratios
        \item Máxima verosimilitud
        \item Extensión a clasificación multiclase: one-vs-rest, softmax
    \end{itemize}
    \item Validación de modelos:
    \begin{itemize}
        \item División de datos: train, validation, test
        \item Validación cruzada: k-fold, leave-one-out, stratified k-fold
        \item Curvas de aprendizaje
    \end{itemize}
\end{itemize}

\subsection{Práctica}
\begin{itemize}[leftmargin=*, labelindent=0pt, label={\textcolor{practicacolor}{$\blacktriangleright$}}]
    \item \textbf{Práctica 3.1:} Regresión lineal simple: ajuste, interpretación, visualización
    \item \textbf{Práctica 3.2:} Análisis de residuos y diagnóstico del modelo
    \item \textbf{Práctica 3.3:} Regresión lineal múltiple: selección de variables
    \item \textbf{Práctica 3.4:} Detección y tratamiento de multicolinealidad
    \item \textbf{Práctica 3.5:} Regularización: Ridge, LASSO, Elastic Net con scikit-learn y glmnet
    \item \textbf{Práctica 3.6:} Selección de hiperparámetros con cross-validation
    \item \textbf{Práctica 3.7:} Regresión logística para clasificación binaria
    \item \textbf{Práctica 3.8:} Interpretación de odds ratios y probabilidades predichas
    \item \textbf{Práctica 3.9:} Validación cruzada: implementación y comparación de modelos
    \item \textbf{Práctica 3.10:} Proyecto: modelado predictivo de un problema real (regresión o clasificación)
\end{itemize}

\section{Unidad 4: Machine Learning Supervisado}

\subsection{Contenidos teóricos}
\begin{itemize}[leftmargin=*]
    \item Clasificación: conceptos fundamentales
    \begin{itemize}
        \item Diferencias entre regresión y clasificación
        \item Clases balanceadas vs desbalanceadas
        \item Estrategias de evaluación
    \end{itemize}
    \item Árboles de decisión:
    \begin{itemize}
        \item Construcción del árbol: criterios de splitting (Gini, entropía, ganancia de información)
        \item Poda: pre-pruning y post-pruning
        \item Interpretabilidad
        \item Ventajas y limitaciones
    \end{itemize}
    \item Random Forest:
    \begin{itemize}
        \item Ensambles: bagging y bootstrap
        \item Votación mayoritaria
        \item Importancia de variables
        \item Hiperparámetros: n\_estimators, max\_depth, min\_samples\_split
    \end{itemize}
    \item Support Vector Machines (SVM):
    \begin{itemize}
        \item Hiperplano óptimo
        \item Margen y vectores de soporte
        \item Kernels: lineal, polinomial, RBF
        \item Parámetros: C (regularización), gamma
    \end{itemize}
    \item K-Nearest Neighbors (KNN):
    \begin{itemize}
        \item Algoritmo basado en instancias
        \item Selección de k
        \item Distancias: euclidiana, Manhattan, Minkowski
        \item Limitaciones: escalabilidad, curse of dimensionality
    \end{itemize}
    \item Métricas de clasificación:
    \begin{itemize}
        \item Matriz de confusión
        \item Accuracy, precision, recall (sensitivity), specificity
        \item F1-score, F-beta score
        \item Curvas ROC y AUC
        \item Curvas precision-recall
        \item Cohen's kappa
    \end{itemize}
    \item Ingeniería de features:
    \begin{itemize}
        \item Interacciones entre variables
        \item Transformaciones no lineales
        \item Extracción de características de texto y fechas
        \item Reducción de dimensionalidad: PCA, selección basada en modelos
    \end{itemize}
    \item Manejo de clases desbalanceadas:
    \begin{itemize}
        \item Técnicas de resampling: oversampling (SMOTE), undersampling
        \item Pesos de clase
        \item Métricas apropiadas (evitar accuracy)
        \item Estratificación en splits
    \end{itemize}
    \item Comparación y selección de modelos:
    \begin{itemize}
        \item Pipelines de scikit-learn
        \item Grid search y random search
        \item Consideraciones: interpretabilidad, tiempo de entrenamiento, escalabilidad
    \end{itemize}
\end{itemize}

\subsection{Práctica}
\begin{itemize}[leftmargin=*, labelindent=0pt, label={\textcolor{practicacolor}{$\blacktriangleright$}}]
    \item \textbf{Práctica 4.1:} Árboles de decisión: construcción, visualización, interpretación
    \item \textbf{Práctica 4.2:} Random Forest: ajuste y análisis de importancia de variables
    \item \textbf{Práctica 4.3:} SVM: experimentación con diferentes kernels
    \item \textbf{Práctica 4.4:} KNN: selección de k y análisis de distancias
    \item \textbf{Práctica 4.5:} Cálculo de métricas de clasificación: matriz de confusión, precision, recall, F1
    \item \textbf{Práctica 4.6:} Curvas ROC y selección de umbrales
    \item \textbf{Práctica 4.7:} Ingeniería de features: creación y selección de variables
    \item \textbf{Práctica 4.8:} Tratamiento de clases desbalanceadas: SMOTE y pesos de clase
    \item \textbf{Práctica 4.9:} Comparación de modelos: grid search, pipelines
    \item \textbf{Práctica 4.10:} Proyecto: clasificación en un dataset real con clases desbalanceadas
\end{itemize}

\section{Unidad 5: Predicción y Comunicación de Resultados}

\subsection{Contenidos teóricos}
\begin{itemize}[leftmargin=*]
    \item Ciclo de vida de un modelo en producción:
    \begin{itemize}
        \item Entrenamiento, validación, deployment, monitoreo
        \item Concept drift y data drift
        \item Reentrenamiento y versionado de modelos
    \end{itemize}
    \item Serialización de modelos:
    \begin{itemize}
        \item Formatos: pickle, joblib, ONNX
        \item Consideraciones de seguridad y compatibilidad
    \end{itemize}
    \item Interpretabilidad de modelos:
    \begin{itemize}
        \item Modelos interpretables vs caja negra
        \item SHAP (SHapley Additive exPlanations)
        \item LIME (Local Interpretable Model-agnostic Explanations)
        \item Partial dependence plots
        \item Importancia de features
    \end{itemize}
    \item Comunicación de resultados:
    \begin{itemize}
        \item Audiencias técnicas vs no técnicas
        \item Estructura de un reporte: contexto, metodología, resultados, conclusiones
        \item Visualizaciones efectivas para comunicar hallazgos
        \item Limitaciones y advertencias
    \end{itemize}
    \item Reportes reproducibles:
    \begin{itemize}
        \item Jupyter notebooks: narrativa, código, visualizaciones
        \item R Markdown: integración de código R, LaTeX, markdown
        \item Exportación a HTML, PDF, presentaciones
        \item Documentación de código y análisis
    \end{itemize}
    \item Ética en análisis de datos:
    \begin{itemize}
        \item Sesgo en datos y modelos
        \item Fairness y discriminación algorítmica
        \item Privacidad y protección de datos
        \item Responsabilidad en el uso de modelos predictivos
    \end{itemize}
\end{itemize}

\subsection{Práctica}
\begin{itemize}[leftmargin=*, labelindent=0pt, label={\textcolor{practicacolor}{$\blacktriangleright$}}]
    \item \textbf{Práctica 5.1:} Serialización y carga de modelos: pickle, joblib
    \item \textbf{Práctica 5.2:} Interpretabilidad: SHAP values y feature importance
    \item \textbf{Práctica 5.3:} LIME: explicaciones locales de predicciones
    \item \textbf{Práctica 5.4:} Creación de reportes con Jupyter notebooks
    \item \textbf{Práctica 5.5:} Reportes con R Markdown: integración de análisis y narrativa
    \item \textbf{Práctica 5.6:} \textbf{Proyecto integrador final:} Análisis completo de un problema real
    \begin{itemize}
        \item Obtención y limpieza de datos
        \item Exploración y visualización
        \item Modelado y evaluación
        \item Comunicación de resultados en formato de reporte profesional
    \end{itemize}
\end{itemize}

\section*{Evaluación}

\begin{itemize}[leftmargin=*]
    \item \textbf{Prácticas semanales:} 40\% (entrega de notebooks resueltos)
    \item \textbf{Parcial:} 30\% (evaluación práctica en computadora)
    \item \textbf{Proyecto final integrador:} 30\% (análisis completo + reporte + presentación)
\end{itemize}

\section*{Bibliografía}

\subsection*{Sugerida}
\begin{itemize}[leftmargin=*]
    \item Wickham, H., \& Grolemund, G. (2017). \textit{R for Data Science}. O'Reilly Media.
    \item VanderPlas, J. (2016). \textit{Python Data Science Handbook}. O'Reilly Media.
    \item James, G., Witten, D., Hastie, T., \& Tibshirani, R. (2021). \textit{An Introduction to Statistical Learning with Applications in R} (2nd ed.). Springer.
\end{itemize}

\subsection*{Complementaria}
\begin{itemize}[leftmargin=*]
    \item McKinney, W. (2022). \textit{Python for Data Analysis} (3rd ed.). O'Reilly Media.
    \item Géron, A. (2022). \textit{Hands-On Machine Learning with Scikit-Learn, Keras, and TensorFlow} (3rd ed.). O'Reilly Media.
    \item Wickham, H. (2016). \textit{ggplot2: Elegant Graphics for Data Analysis} (2nd ed.). Springer.
    \item Kuhn, M., \& Johnson, K. (2019). \textit{Feature Engineering and Selection: A Practical Approach for Predictive Models}. Chapman and Hall/CRC.
\end{itemize}

\end{document}
